\section{Resultados y análisis}
\subsection{Datos experimentales}
\textit{Pendiente: presentar las mediciones en tablas con unidades e incertidumbres disponibles. Conservar los datos originales en data/.}

% Las tablas pueden guardarse en tables/ e incluirse con:
% \begin{table}[htbp]
\centering
\caption{Datos experimentales registrados}
\label{tab:datos}
\begin{tabular}{ccc}
\toprule
\textbf{Medición} & \textbf{$x$ (\si{unidad})} & \textbf{$y$ (\si{unidad})} \\
\midrule
1 & $0.00 \pm 0.00$ & $0.00 \pm 0.00$ \\
2 & $0.00 \pm 0.00$ & $0.00 \pm 0.00$ \\
3 & $0.00 \pm 0.00$ & $0.00 \pm 0.00$ \\
\bottomrule
\end{tabular}
\end{table}

% Usar booktabs y ajustar su ancho a \columnwidth.

\subsection{Cálculos y gráficas}
\textit{Pendiente: desarrollar los cálculos e incorporar las gráficas y los ajustes necesarios para responder los objetivos.}

% Guardar gráficas exportadas en images/ y figuras TikZ/PGFPlots en figures/.
% Ejemplo de inclusión de una figura vectorial:
% \input{figures/ajuste}

\subsection{Comparación con el modelo teórico}
\textit{Pendiente: comparar los resultados experimentales con las predicciones del modelo y cuantificar las discrepancias cuando corresponda.}
